\section{Orientation covers and descent of isotropic families}
\label{sec:orientation-descent}
We give the descent statements used in the component theorem explicitly.
They distinguish connectedness of a cover from geometric irreducibility
of a fibre, and also distinguish one maximal-isotropic factor from a
product of such factors.

\begin{proposition}[Relative orthogonal descent]\label{prop:orientation-descent}
Let $S$ be a smooth integral complex variety and let $(\mathcal V,q)$
be a nondegenerate quadratic bundle of rank $2s$, with values in a line
bundle. Its orientation cover $\widetilde S\to S$ is the degree-two
etale cover whose local equation, after choosing frames, is
$z^2=(-1)^s\det q$. Changes of frames glue these local covers.
Suppose this cover is connected. The relative maximal orthogonal
Grassmannian $\operatorname{OGr}_S(s,\mathcal V)$ is integral.
More generally, the product of $L$ such Grassmannians for the same
quadratic bundle has $2^{L-1}$ irreducible components, indexed by
$\{+,-\}^L$ modulo simultaneous sign change.
Each component is smooth over $S$. Adding relative nonmaximal
orthogonal Grassmannian factors does not change this count.
\end{proposition}
\begin{proof}
If a frame of $\mathcal V$ changes by $A$, its determinant changes
by $(\det A)^2$. If the frame of the value line changes by $u$,
the determinant changes by $u^{-2s}$. Thus the square-root equations
are compatible with the transition functions of
$\det\mathcal V^*\otimes(\text{value line})^s$. Their zero-free
right-hand side makes the cover etale. After an etale base change the
form has a hyperbolic basis $e_1,\ldots,e_s,f_1,\ldots,f_s$.

Over an algebraically closed field a maximal isotropic plane belongs
to one of two families, distinguished by the parity of its intersection
dimension with $\Span(e_1,\ldots,e_s)$. A Witt basis sends any such
plane to the reference plane by an orthogonal transformation. The
transformations of determinant one preserve this parity and act
transitively on each family; interchanging $e_1$ and $f_1$ reverses it.
The special orthogonal group is connected: in a hyperbolic basis its
elementary orthogonal root subgroups are affine lines, its diagonal
torus is a product of multiplicative groups, and these generate the
group by the usual elimination of a Witt basis. Therefore each of its
two orbits is irreducible. Equivalently, the graph charts within a
family are alternating matrices and are translates under this
connected group. The same Witt-basis elimination gives a single
irreducible orbit for isotropic planes of dimension strictly less
than half the rank.

The two families are consequently locally constant etale labels,
with the determinant character describing their interchange. The
orientation cover kills this character. Over $\widetilde S$ each
family is a smooth proper relative variety with geometrically
irreducible fibres. It is integral: smoothness gives reducedness,
every component dominates the integral base, and the geometric
generic fibre has only one component. Here domination follows also
from local smoothness and properness: each component is open and
closed in the smooth total space, so its nonempty image is open and
closed in the connected base. The nontrivial deck transformation
exchanges the two families. Their images in the descended space are
therefore the same irreducible image, which is the whole space.
This proves integrality for one factor.

For $L$ factors the geometric components over $\widetilde S$ are the
$2^L$ products of these integral family bundles. The deck action sends
$\varepsilon$ to $-\varepsilon$. The orbit unions descend to the
irreducible components: the image of any one integral product is
irreducible and equals the image of its entire orbit, and distinct
orbits have distinct geometric generic fibres. There are $2^{L-1}$
orbits. Smoothness descends etale locally. Nonmaximal factors have
geometrically irreducible fibres and add no labels.
\end{proof}

\begin{proposition}[The five exceptional orientation bases]
\label{prop:exceptional-orientation-base}
For each exceptional tuple of Theorem~\ref{thm:all-dimension}, the
incidence containing the full radical and a maximal isotropic quotient
is irreducible. For the full-rank case $(c,k,L)=(7,14,1)$ the
nondegenerate isotropic incidence is irreducible as well.
\end{proposition}
\begin{proof}
For an exceptional rank $r=2s<k$, work first over the dense ordered
support chart
\[
 T_r=\{(x_1,\ldots,x_r,\alpha_1,\ldots,\alpha_{r-1}):
 x_i\ne x_j,\ \alpha_i\ne0\},\qquad \alpha_r=1.
\]
This is a smooth integral open subset of
$\mathbb A^r\times\mathbb G_m^{r-1}$. It parameterizes the
projective functionals $[\ell]=[\sum_i\alpha_i\operatorname{ev}_{x_i}]$.
Evaluation identifies the quotient by the radical with $\C^r$,
and the induced quadratic form is $\diag(\alpha_1,\ldots,\alpha_r)$.
In the function field
$\C(x_1,\ldots,x_r,\alpha_1,\ldots,\alpha_{r-1})$, the product
$\prod_{i<r}\alpha_i$ is not a square: its valuation at $\alpha_1=0$
is one. This valuation is a function-field valuation; the boundary
$\alpha_1=0$ need not belong to $T_r$. The orientation cover obtained
by adjoining its square root is therefore integral and connected.
Proposition~\ref{prop:orientation-descent} applies.

The radical is $gV_{k-r-1}$ with $g=\prod_i(t-x_i)$. Containing
this radical identifies the exceptional plane incidence with the
maximal orthogonal Grassmannian of the quotient. There is one such
factor in each of the five cases. The ordered chart maps dominantly
to the corresponding exact-rank base. Its generic fibres merely
relabel the support: the radical determines $g$, its distinct roots
recover the unordered points, and a Vandermonde system recovers the
weights up to common scale. Taking the unordered quotient cannot
split the irreducible image of this incidence. Its closure is the
exceptional incidence claimed in the theorem. Colliding-support
points do not create another component of this closure.

In the remaining full-rank case the universal quadratic form has
even rank fourteen. The determinant of the universal Hankel matrix
has the irreducible corank-one Hankel hypersurface as its zero divisor,
with multiplicity one. The irreducibility follows from the Hankel
maximal-minor description in Section~\ref{sec:global-geometry}; simplicity can
also be checked at a corank-one form represented by thirteen distinct
evaluations. Its kernel is generated by a nonzero polynomial $v$;
the derivative of the determinant in a Hankel direction $\dot\ell$
is a nonzero scalar multiple of $\dot\ell(v^2)$, which is not
identically zero. On an affine projective chart meeting its generic
point this simple prime divisor proves nonsquareness in the function
field. Scaling the representative of a rank-fourteen form changes
the determinant by a fourteenth power, hence by a square, so this is
the projective orientation cover. It is connected and
Proposition~\ref{prop:orientation-descent} gives the assertion.
\end{proof}
